\documentclass[10pt,a4paper]{amsart}
\usepackage{amsmath,amssymb,amsthm,graphicx}
\usepackage{hyperref}
\numberwithin{equation}{section}

\newtheorem{Theorem}{Theorem}[section]
\newtheorem*{Theorem*}{Theorem}
\newtheorem{Corollary}[Theorem]{Corollary}
\newtheorem{Lemma}[Theorem]{Lemma}
\newtheorem{Proposition}[Theorem]{Proposition}

\theoremstyle{definition}
\newtheorem{Definition}[Theorem]{Definition}
\newtheorem{Note}[Theorem]{Note}
\newtheorem{Example}[Theorem]{Example}
\newtheorem{Remark}[Theorem]{Remark}
\DeclareMathOperator{\Hom}{Hom}
\DeclareMathOperator{\End}{End}
\DeclareMathOperator{\Ker}{Ker}
\DeclareMathOperator{\re}{re}
\DeclareMathOperator{\Ima}{Im}
\DeclareMathOperator{\ev}{ev}
\DeclareMathOperator{\id}{id}
\DeclareMathOperator{\tr}{tr}
\DeclareMathOperator{\ad}{ad}
\DeclareMathOperator{\mult}{mult}
\DeclareMathOperator{\fin}{fin}
\DeclareMathOperator{\str}{str}
\DeclareMathOperator{\row}{row}
\DeclareMathOperator{\col}{col}
\newcommand{\ilim}[1][]{\mathop{\varinjlim}\limits_{#1}}
\newcommand{\plim}[1][]{\mathop{\varprojlim}\limits_{#1}}
\newcommand{\ve}{\varepsilon}
\newcommand{\ts}{\,}
\newcommand{\tss}{\hspace{1pt}}
\newcommand{\on}{\operatorname}
\newcommand{\relmiddle}[1]{\mathrel{}\middle#1\mathrel{}}
\renewcommand{\labelenumii}{\arabic{enumii}).}
\allowdisplaybreaks
\title{Injectivity of the omega-fixed central-extension presentation inside the universal central extension of sl(n)[u\textasciicircum{}\{+-1\},v]}
\author{Mamoru Ueda}
\date{doi:10.3842/SIGMA.2026.089; CC BY 4.0}
\begin{document}
\maketitle
\noindent\textit{Allowed standard prerequisites (editorial).}\begin{itemize}
\item PBW theorem, affine root systems, universal central extensions. Original MRY, GNW and HU bibliography entries; author formulations and uses preserved
\end{itemize}\medskip
\section[Universal central extension of sl(n){[u\^{}\{pm1\},v]}]{Universal central extension of $\boldsymbol{\mathfrak{sl}(n)\bigl[u^{\pm1},v\bigr]}$}
In this section, we recall the work of~\cite{MRY}, which gives a presentation of the universal central extension of $\mathfrak{sl}(n)\bigl[u^{\pm1},v\bigr]$. For a perfect Lie algebra $\mathfrak{g}$, a central extension of $\mathfrak{g}$ is a pair consisting of a~Lie algebra \smash{$\widetilde{\mathfrak{g}}$} and a surjective homomorphism $\pi\colon\widetilde{\mathfrak{g}}\to\mathfrak{g}$ whose kernel is contained in the center of~\smash{$\widetilde{\mathfrak{g}}$}. The central extension $(\widetilde{\mathfrak{g}},\pi)$ is called the universal central extension of $\mathfrak{g}$ if for every central extension $(\mathfrak{e},\psi)$ there exists a unique homomorphism $\phi\colon\widetilde{\mathfrak{g}}\to\mathfrak{e}$ satisfying that $\psi\circ\phi=\pi$.

In~\cite{Kassel}, Kassel gave an explicit presentation of the universal central extension of $\mathfrak{sl}(n)\bigl[u^{\pm1},v\bigr]$.
We set $\bigl(\Omega\bigl(\mathbb{C}\bigl[u^{\pm1},v\bigr]\bigr),d\bigr)$ as the module of differentials of $\mathbb{C}\bigl[u^{\pm1},v\bigr]$ in the same way as in~\cite[Section~2]{MRY}. We take $\{a_i\}$ as a basis of $\mathbb{C}\bigl[u^{\pm1},v\bigr]$. Let $F$ be a free module of $\mathbb{C}\bigl[u^{\pm1},v\bigr]$ whose basis is \smash{$\bigl\{\widetilde{d} a_i\bigr\}$}. We regard $F$ as a 2-sided \smash{$\mathbb{C}\bigl[u^{\pm1},v\bigr]$} module by setting \smash{$a\bigl(\widetilde{d} b\bigr)=\bigl(\widetilde{d} b\bigr)a$} for~any~\smash{$a,b\in\mathbb{C}\bigl[u^{\pm1},v\bigr]$}. We set a linear map \smash{$\widetilde{d} \colon \mathbb{C}\bigl[u^{\pm1},v\bigr]\to F$} as \smash{$\widetilde{d} \bigl(\sum_{i}c_ia_i\bigr)=\sum_{i}c_i\widetilde{d} a_i$} and $K$ as a~submodule of~$F$ generated by \smash{$\widetilde{d}(ab)=(\widetilde{d}a)b+a\bigl(\widetilde{d} b\bigr)$} for $a,b\in\mathbb{C}\bigl[u^{\pm1},v\bigr]$. We denote $F/K$ by~\smash{$\Omega\bigl(\mathbb{C}\bigl[u^{\pm1},v\bigr]\bigr)$}. We also set \smash{$d\colon\mathbb{C}\bigl[u^{\pm1},v\bigr]\to\Omega\bigl(\mathbb{C}\bigl[u^{\pm1},v\bigr]\bigr)$} as a homomorphism induced by \smash{$\widetilde{d}$}.

Let \smash{$\mathfrak{sl}(n)\otimes\mathbb{C}\bigl[u^{\pm1},v\bigr]\oplus\Omega\bigl(\mathbb{C}\bigl[u^{\pm1},v\bigr]\bigr)/d\mathbb{C}\bigl[u^{\pm1},v\bigr]$} be the Lie algebra whose commutator relations are given by
\begin{gather*}
[z_1\otimes a_1,z_2\otimes a_2]=[z_1,z_2]\otimes a_1\cdot a_2+\kappa(z_1,z_2)a_2da_1,\\
\bigl[\mathfrak{sl}(n)\otimes\mathbb{C}\bigl[u^{\pm1},v\bigr],\Omega^1\bigl(\mathbb{C}\bigl[u^{\pm1},v\bigr]\bigr)/d\mathbb{C}\bigl[u^{\pm1},v\bigr]\bigr]=0,\\
\bigl[\Omega^1\bigl(\mathbb{C}\bigl[u^{\pm1},v\bigr]\bigr)/d\mathbb{C}\bigl[u^{\pm1},v\bigr],\Omega^1\bigl(\mathbb{C}\bigl[u^{\pm1},v\bigr]\bigr)/d\mathbb{C}\bigl[u^{\pm1},v\bigr]\bigr]=0.
\end{gather*}
We can naturally define a surjective homomorphism
\[
\pi\colon\ \mathfrak{sl}(n)\otimes\mathbb{C}\bigl[u^{\pm1},v\bigr]\oplus\Omega\bigl(\mathbb{C}\bigl[u^{\pm1},v\bigr]\bigr) / d\mathbb{C}\bigl[u^{\pm1},v\bigr]\to\mathfrak{sl}(n)\otimes\mathbb{C}\bigl[u^{\pm1},v\bigr].
\]
Since the kernel of $\pi$ is $\Omega\bigl(\mathbb{C}\bigl[u^{\pm1},v\bigr]\bigr)/d\mathbb{C}\bigl[u^{\pm1},v\bigr]$, $\pi$ gives a central extension.
Moreover, by~\cite{Kassel}, $\mathfrak{sl}(n)\otimes\mathbb{C}\bigl[u^{\pm1},v\bigr]\oplus\Omega \bigl(\mathbb{C}\bigl[u^{\pm1},v\bigr]\bigr)/d\mathbb{C}\bigl[u^{\pm1},v\bigr]$ together with $\pi$ form the universal central extension of $\mathfrak{sl}(n)\otimes\mathbb{C}\bigl[u^{\pm1},v\bigr]$.

Moody--Rao--Yokonuma~\cite{MRY} gave a new presentation of this universal central extension, which we describe below.
Let us take a matrix \smash{$\bigl(a^A_{i,j}\bigr)_{0\leq i,j\leq n-1}$} as
\begin{equation*}
a^A_{i,j}=\begin{cases}
\hphantom{-}2&\text{if }i=j,\\
-1&\text{if }i=j\pm1,\\
-1&\text{if }(i,j)=(0,n-1),(n-1,0),\\
\hphantom{-}0&\text{otherwise}.
\end{cases}
\end{equation*}
Then, \smash{$\bigl(a^A_{i,j}\bigr)_{0\leq i,j\leq n-1}$} is the Cartan matrix of an affine Lie algebra $\smash{\widehat{\mathfrak{sl}}(n)}=\mathfrak{sl}(n)\otimes\mathbb{C}\bigl[t^{\pm1}\bigr]\oplus\mathbb{C}c$ whose commutator relations are given by
\begin{gather*}
[x\otimes t^u,y\otimes t^s]=[x,y]\otimes t^{u+s}+u\delta_{u+s,0}\tr(xy)c \qquad\text{for }x,y\in\mathfrak{sl}(n),\\
\bigl[c,\widehat{\mathfrak{sl}}(n)\bigr]=0,
\end{gather*}
where $\tr$ is the trace of $\mathfrak{gl}(n)$.
\begin{Definition}%\label{Prop32}
We define $\mathfrak{A}(n)$ as the Lie algebra over $\mathbb{C}$ generated by
\begin{equation*}
\bigl\{x_{i,r}^\pm, h_{i,r}\mid 0\leq i\leq n-1,\,r\in\mathbb{Z}_{\geq0}\bigr\}
\end{equation*}
subject to the following defining relations:
\begin{gather}
[h_{i,r}, h_{j,s}] = 0,\label{Eq2.1}\\
\bigl[x_{i,r}^{+}, x_{j,s}^{-}\bigr] = \delta_{i,j} h_{i, r+s},\label{Eq2.2}\\
\bigl[h_{i,0}, x_{j,r}^{\pm}\bigr] = \pm a^A_{i,j} x_{j,r}^{\pm},\label{Eq2.4}\\
\bigl[h_{i,r+1}, x_{j,s}^{\pm}\bigr] = \bigl[h_{i,r}, x_{j,s+1}^{\pm}\bigr],\label{Eq2.5}\\
\bigl[x_{i, r+1}^{\pm}, x_{j, s}^{\pm}\bigr] = \bigl[x_{i, r}^{\pm}, x_{j, s+1}^{\pm}\bigr],\label{Eq2.8}\\
\sum_{\sigma\in S_{1-a_{i,j}}}\limits\bigl(\ad x_{i,r_{\sigma(1)}}^{\pm}\bigr)\cdots\bigl(\ad x_{i,r_{\sigma(1-a_{i,j})}}^{\pm}\bigr)\bigl(x_{j,s}^{\pm}\bigr)= 0 \qquad\text{if }i \neq j,\label{Eq2.10}
\end{gather}
where $S_l$ is a symmetric group of degree $l$. We also denote the universal enveloping algebra of~$\mathfrak{A}(n)$~by~$A(n)$.
\end{Definition}
By using \eqref{Eq2.8}, we obtain
\begin{equation*}
\bigl[x^\pm_{i,r},x^\pm_{i,s}\bigr]=\frac{1}{2}\bigl(\bigl[x^\pm_{i,r+s},x^\pm_{i,0}\bigr]+\bigl[x^\pm_{i,0},x^\pm_{i,r+s}\bigr]\bigr)=0.%\label{Eq2.11}
\end{equation*}

Moody--Rao--Yokonuma showed that $\mathfrak{A}(n)$ is isomorphic to the universal central extension of~$\mathfrak{sl}(n)\otimes\mathbb{C}\bigl[u^{\pm1},v\bigr]$.
\begin{Theorem}[{\cite[Proposition 3.5]{MRY}}]
There exists an isomorphism
\begin{equation*}
\iota\colon\ \mathfrak{A}(n)\to \mathfrak{sl}(n)\otimes\mathbb{C}\bigl[u^{\pm1},v\bigr]\oplus\Omega\bigl(\mathbb{C}\bigl[u^{\pm1},v\bigr]\bigr)/d\mathbb{C}\bigl[u^{\pm1},v\bigr]
\end{equation*}
given by
\begin{gather*}
x^+_{i,r}=\begin{cases}
E_{i,i+1}v^r&\text{if }i\neq0,\\
E_{n,1}uv^r&\text{if }i=0,
\end{cases}\qquad
x^-_{i,r}=\begin{cases}
E_{i+1,i}v^r&\text{if }i\neq0,\\
E_{1,n}u^{-1}v^r&\text{if }i=0,
\end{cases}
\end{gather*}
where $E_{i,j}$ is a matrix unit of $\mathfrak{gl}(n)$ whose $(p,q)$ component is $\delta_{p,i}\delta_{q,j}$.
\end{Theorem}

Finally, we consider a presentation of the subalgebra of $A(n)$ generated by $\bigl\{x^+_{i,r}\bigr\}$.
We take a~subspace \smash{$\Omega_+\subset\Omega\bigl(\mathbb{C}\bigl[u^{\pm1},v\bigr]\bigr)/d\mathbb{C}\bigl[u^{\pm1},v\bigr]$} spanned by \smash{$\bigl\{u^{r-1}v^sdu\mid r\geq1,\,s\geq0\bigr\}$}.
We set a~subalgebra $A(n)_+\subset A(n)$ generated by \smash{$\bigoplus_{s\geq0}\bigoplus_{\alpha\in\Delta^+}\mathfrak{g}_\alpha v^s$} and $\Omega_+$. Let us take a grading on~$A(n)$ by \smash{${\rm deg}(h_{i,r})={\rm deg}\bigl(x^\pm_{i,r}\bigr)=r+1$} and denote the graded algebra associated with~$A(n)$~by~$\operatorname{gr}A(n)$.
\begin{Theorem}\label{plus}
The associative algebra $A(n)_+$ is isomorphic to the associative algebra generated by $\bigl\{\bar{x}^+_{i,r}\mid 1\leq i\leq n-1,\,r\geq0\bigr\}$ with the relations:
\begin{gather*}
\bigl[\bar{x}^+_{i,r+1},\bar{x}^+_{j,s}\bigr]=\bigl[\bar{x}^+_{i,r},\bar{x}^+_{j,s+1}\bigr],\\
\bigl[\bar{x}^+_{i,r},\bar{x}^+_{j,s}\bigr]=0 \qquad \text{if }a^A_{i,j}=0,\\
\bigl[\bar{x}^+_{i,r},\bigl[\bar{x}^+_{i,s},\bar{x}^+_{j,u}\bigr]\bigr]=0 \qquad \text{if }a^A_{i,j}=-1.
\end{gather*}
\end{Theorem} % RS: pre -> {\rm pre}
Let $\mathfrak{A}(n)_{{\rm pre}}$ be the Lie algebra generated by \smash{$\bigl\{h_{i,r},x^\pm_{i,r}\mid 1\leq i\leq n-1,\,r\geq0\bigr\}$} with the defining relations \eqref{Eq2.1}--\eqref{Eq2.8}. We also define $\mathfrak{A}(n)_{{\rm pre},\pm}$ and $\mathfrak{A}(n)_{{\rm pre},0}$ as Lie subalgebras generated by \smash{$\bigl\{x^\pm_{i,r},h_{i,r}\mid 1\leq i\leq n-1,\,r\geq0\bigr\}$}.
Theorem~\ref{plus} follows from the following lemma.
\begin{Lemma}\qquad
\begin{enumerate}\itemsep=0pt
\item[$(1)$] The associative algebra \raisebox{0.3pt}{$U(\mathfrak{A}(n)_{{\rm pre},\pm})$} is isomorphic to a quotient algebra of the free associative algebra generated by \smash{\raisebox{-0.3pt}{$\bigl\{x^\pm_{i,r}\bigr\}$}} divided by \eqref{Eq2.8}.

\item[$(2)$] Let $I_\pm$ be an ideal of $\mathfrak{A}(n)_{{\rm pre}}$ generated by \eqref{Eq2.10}. Then, $I_-\cap A(n)_{{\rm pre},+}=\varnothing$.
\end{enumerate}
\end{Lemma}
\begin{proof} % RS: не розумію що це, але запис трохи дивно виглядає
\textup{(1)}\ Let $V$ be a quotient algebra of a noncommutative polynomial ring $\mathbb{C}\langle y_{i,d}\mid0\leq i\leq n-1,d\in\mathbb{Z}_{\geq0}\rangle$ divided by the relations $[y_{i,r+1},y_{j,s}]-[y_{i,r},y_{j,s+1}]$. By a direct computation, we obtain a homomorphism
\begin{align*}
\kappa\colon\ \mathfrak{A}(n)_{{\rm pre}}\to\mathfrak{gl}(V)
\end{align*}
by
\begin{align*}
\begin{split}
&\bigl(\kappa\bigl(x^-_{i,d}\bigr)\bigr)(y_{i_1,d_1}\cdots y_{i_r,d_r})=y_{i,d}y_{i_1,d_1}\cdots y_{i_r,d_r},\\
&(\kappa(h_{i,d}))(y_{i_1,d_1}\cdots y_{i_r,d_r})=-\sum_{u=1}^r\limits a^A_{i,i_u}y_{i_1,d_1}\cdots y_{i_{u-1},d_{u-1}}y_{i_u,d+d_u}y_{i_{u+1},d_{u+1}}\cdots y_{i_r,d_r},\\
&\bigl(\kappa\bigl(x^+_{i,d}\bigr)\bigr)(1)=0,\\
&\bigl(\kappa\bigl(x^+_{i,d}\bigr)\bigr)(y_{i_1,d_1}\cdots y_{i_r,d_r})=\delta_{i,i_1}\kappa(h_{i,d+d_1})(y_{i_2,d_2}\cdots y_{i_r,d_r})+y_{i_1,d_1}\bigl(\kappa\bigl(x^+_{i,d}\bigr)\bigr)(y_{i_2,d_2}\cdots y_{i_r,d_r}).
\end{split}
\end{align*}
The homomorphism $\kappa$ induces a surjective homomorphism $\Xi\colon \mathfrak{A}(n)_{{\rm pre},-}\to V$ by $\Xi(x)=(\kappa(x))1$.
By a direct computation, we can also construct a homomorphism
\begin{align*}
\xi\colon\ V\to \mathfrak{A}(n)_{{\rm pre},-},\qquad y_{i,d}\mapsto x^-_{i,d}.
\end{align*}
Since $\xi\circ\Xi$ is an identity map, $\Xi$ is bijective. Then, we find that $U(\mathfrak{A}(n)_{{\rm pre},-})$ is isomorphic to a quotient algebra of the free associative algebra generated by $\{x^-_{i,r}\}$ divided by \eqref{Eq2.8}. Since~$\mathfrak{A}(n)_{{\rm pre},+}$ is isomorphic to $\mathfrak{A}(n)_{{\rm pre},-}$, $U(\mathfrak{A}(n)_{{\rm pre},+})$ is also isomorphic to a quotient algebra of the free associative algebra generated by $\bigl\{x^+_{i,r}\mid r\geq0\bigr\}$ divided by \eqref{Eq2.8}.

\textup{(2)}\ It is enough to show that $I_-$ is contained in $\mathfrak{A}(n)_{{\rm pre},-}$. This follows from the fact that $\ad\bigl(x^+_{i,r}\bigr)$ acts on the left-hand side of \eqref{Eq2.10} trivially.
\end{proof}

\section[A central extension of sl(n){[u\^{}\{pm1\},v]}\^{}omega]{A central extension of $\boldsymbol{\mathfrak{sl}(n)\bigl[u^{\pm1},v\bigr]^\omega}$}\label{sec:3}
First, we set notations about the affine Lie algebra $\widehat{\mathfrak{sl}}(n)$.
Let $\Delta$ \smash{\big(resp.\ $\Delta^+$, $\Delta^-$\big)} be the set of roots (resp.\ positive roots, negative roots) of $\widehat{\mathfrak{sl}}(n)=\mathfrak{sl}(n)\bigl[u^{\pm1}\bigr]\oplus\mathbb{C}c$.
We set an inner product~$(\, ,\, )$ on~\smash{$\widehat{\mathfrak{sl}}(n)$} by
\begin{equation*}
(xu^s,yu^r)=\tr(xy)\delta_{r+s,0},\qquad (xu^r,c)=(c,c)=0.
\end{equation*}
We denote by $\mathfrak{g}_\alpha$ the root space associated with a root $\alpha$.
For $\alpha\in\Delta$, we take \smash{$\bigl\{x_\alpha^k\bigr\}_{k=1}^{\dim (\mathfrak{g}_\alpha)}$} as a~basis of $\mathfrak{g}_\alpha$ such that \smash{$\bigl(x_\alpha^k,x_\beta^l\bigr)=\delta_{\alpha+\beta,0}\delta_{k,l}$} and \smash{$\bigl(x_\alpha^k\bigr)^{\mathsf{T}}=x_{-\alpha}^k$}, where $X^{\mathsf{T}}$ is the transpose of~a~matrix~$X$.
The following relations hold in \smash{$U\bigl(\widehat{\mathfrak{sl}}(n)\bigr)^{\otimes2}$} (see~\cite[Corollary 2.4]{K}):
\begin{gather}
\sum_{k=1}^{\dim \mathfrak{g}_\alpha}\bigl[x_{-\alpha}^k,x_\beta\bigr]\otimes x_\alpha^k=\sum_{k=1}^{\dim \mathfrak{g}_{\beta-\alpha}}x_{\beta-\alpha}^k\otimes \bigl[x_\beta,x_{\alpha-\beta}^k\bigr].\label{out}
\end{gather}
Let $\widehat{\tau}$ be an automorphism of \smash{$U\bigl(\widehat{\mathfrak{sl}}(n)\bigr)$} defined by $X\mapsto-X^{\mathsf{T}}$. Considering the images of both sides of \eqref{out} via $\widehat{\tau}\otimes\id$, we obtain
\begin{gather}
\sum_{k=1}^{\dim \mathfrak{g}_\alpha}\bigl[x_{\alpha}^k,x_{-\beta}\bigr]\otimes x_\alpha^k=-\sum_{k=1}^{\dim \mathfrak{g}_{\beta-\alpha}}x_{\alpha-\beta}^k\otimes \bigl[x_\beta,x_{\alpha-\beta}^k\bigr].\label{out1}
\end{gather}

Let us take an involution of $\mathfrak{sl}(n)\bigl[u^{\pm1},v\bigr]$:
\begin{equation*}
\omega\colon\ \mathfrak{sl}(n)\bigl[u^{\pm1},v\bigr]\to\mathfrak{sl}(n)\bigl[u^{\pm1},v\bigr],\qquad Xu^rv^s\mapsto -(-1)^{s}X^{\mathsf{T}}u^{-r}v^{s}.
\end{equation*}
In this section, we give a central extension of $\mathfrak{sl}(n)\bigl[u^{\pm1},v\bigr]^\omega$.
\begin{Definition}
Let $\mathfrak{BD}(n)$ be the Lie algebra generated by
\begin{gather*}
\{H_{i,2r+1},B_{i,r}\mid0\leq i\leq n-1,\, r\geq0\}
\end{gather*}
with the following relations:
\begin{gather*}
[H_{i,2r+1},H_{j,2s+1}]=0,\\
[H_{i,2r+1},B_{j,s}]=2a^A_{i,j}B_{j,2r+s+1},\\
[B_{i,r+1},B_{j,s}]-[B_{i,r},B_{j,s+1}]=-\delta_{i,j}((-1)^r+(-1)^s)H_{i,r+s},\\
[B_{i,r},B_{j,s}]=0\qquad \text{if }a^A_{i,j}=0,\\
[B_{i,r},[B_{i,s},B_{j,u}]]=-B_{j,r+s+u}\qquad \text{if }a^A_{i,j}=-1.
\end{gather*}
We also denote the universal enveloping algebra of $\mathfrak{BD}(n)$ by $BD(n)$.
\end{Definition}
In~\cite{Gu1}, it was shown that the affine Yangian associated with \smash{$\widehat{\mathfrak{sl}}(n)$} coincides with $A(n)$ if all parameters are zero. If the twisted affine Yangian of types $BD$ exists, we expect that the twisted affine Yangian of types $BD$ should be equal to $BD(n)$ if all the parameters are zero. Moreover, since the twisted affine Yangian should be a coideal of the affine Yangian, it is expected that~$BD(n)$ can be embedded into $A(n)$.

By a direct computation, we obtain a homomorphism $\phi\colon BD(n)\to A(n)$ given by
\begin{align*}
&B_{i,r}=\begin{cases}
E_{i,i+1}v^r-(-1)^rE_{i+1,i}v^r&\text{if }1\leq i\leq 2n-1,\\
E_{n,1}uv^r-(-1)^rE_{1,n}u^{-1}v^r&\text{if }i=0,
\end{cases}\\
&H_{i,2r+1}=\begin{cases}
2(E_{i,i}-E_{i+1,i+1})v^{2r+1}&\text{if }1\leq i\leq 2n-1,\\
2(E_{n,n}-E_{1,1})v^{2r+1}+2u^{-1}v^{2r+1}du&\text{if }i=0.
\end{cases}
\end{align*}

\begin{Theorem}
The homomorphism $\phi$ is injective.
\end{Theorem}

\begin{proof}

For $(i,j,r,s)$ such that $E_{i,j}u^rv^s\in A(n)_+$, we define $E_{i,j}^{r,s}\in BD(n)$ inductively as follows:
\begin{gather*}
E^{0,s}_{i,i+1}=B_{i,s},E^{1,s}_{n,1}=B_{0,s}
\end{gather*}
and
\begin{alignat*}{3}
&E_{i,j}^{r,s}=\bigl[E_{i,j-1}^{r,s},B_{j-1,0}\bigr]\qquad&& \text{if }j>i+1\text{ and }1\leq i<j\leq n,&\\
&E_{i,i+1}^{r,s}=\frac{1}{2}\bigl[H^{r}_{i,s},B_{i,0}\bigr],&\\
&H^r_{i,s}=\bigl[E^{r-1,0}_{i,i+1},E_{i+1,i}^{1,s}\bigr]\qquad&& \text{if }1\leq i\leq n-1,&\\
&H^r_{0,s}=\bigl[B_{0,0},E_{1,n}^{r-1,s}\bigr],&\\
&E_{i,j}^{r,s}=\bigl[E_{i,j-1}^{r,s},B_{j-1,0}\bigr]\qquad&& \text{if }i>j\text{ and }n\geq i\geq j\geq2,&\\
&E_{i,1}^{r,s}=\bigl[E_{i,n}^{r-1,s},B_{0,0}\bigr]\qquad&& \text{if }1\leq i\leq n-1,&\\
&E_{n,1}^{r,s}=\frac{1}{2}\bigl[H^1_{0,0},E_{n,1}^{r-1,s}\bigr]\qquad&&\text{for }r\geq2.&
\end{alignat*}
By a direct computation, we obtain
\begin{alignat*}{3}
&\phi\bigl(E_{i,j}^{r,s}\bigr)=F_{i,j}^{r,s}\qquad&& \text{if }j>i+1\text{ and }1\leq i<j\leq n,&\\
&\phi\bigl(E_{i,i+1}^{r,s}\bigr)=F_{i,i+1}^{r,s},&\\
&\makebox[0pt][l]{$\phi\bigl(H^r_{i,s}\bigr)=F_{i,i}^{r,s}-F_{i+1,i+1}^{r,s}+(r-1)u^{r-1}v^sdu-(r-1)(-1)^su^{-1-r}v^{s}du\qquad\!\! \text{if }1\leq i\leq n-1,$}&\\
&\makebox[0pt][l]{$\phi\bigl(H^r_{0,s}\bigr)=F_{n,n}^{r,s}-F_{1,1}^{r,s}+u^{r-1}v^sdu-(-1)^su^{-1-r}v^{s}du \qquad \text{for }r\geq2,$}&\\
&\phi\bigl(E_{i,j}^{r,s}\bigr)=F_{i,j}^{r,s}\qquad&& \text{if }i>j\text{ and }n\geq i\geq j\geq2,&\\
&\phi\bigl(E_{i,1}^{r,s}\bigr)=F_{i,1}^{r,s}\qquad&& \text{if }1\leq i\leq n-1,&\\
&\phi\bigl(E_{n,1}^{r,s}\bigr)=F_{n,1}^{r,s}\qquad&& \text{for }r\geq2,&
\end{alignat*}
where \smash{$F_{i,j}^{r,s}=E_{i,j}u^rv^s-(-1)^{s}E_{j,i}u^{-r}v^s$}.
Since \smash{$\bigl\{\phi\bigl(E_{i,j}^{r,s}\bigr)\bigr\}\cup\{\phi(H_{i,2r+1})\}$} are linearly independent, the set \smash{$\bigl\{E_{i,j}^{r,s}\bigr\}\cup\{H_{i,2r+1}\}$} is linearly independent.
By the PBW theorem, if we fix an order of the set \smash{$\bigl\{H_{k,2r+1},E_{i,j}^{r,s}\bigr\}$} satisfying that $H_{k,2r+1}<E_{i,j}^{r,s}$, the set of ordered polynomials derived from \smash{$\bigl\{H_{k,2r+1},E_{i,j}^{r,s}\bigr\}$} is linearly independent in $BD(n)$. Thus, it is enough to show that these ordered polynomials span $BD(n)$.

Let us take a grading on $BD(n)$ by ${\rm deg}(B_{j,r})=r+1$ and ${\rm deg}(H_{j,r})=r+1$. We denote the graded algebra associated with $BD(n)$ by $\operatorname{gr}BD(n)$ and the images of \smash{$\bigl\{E^{r,s}_{i,j}\bigr\}$} and \smash{$\bigl\{H^r_{i,s}\bigr\}$} in $\operatorname{gr}BD(n)$ by \smash{$\bigl\{\bar{E}^{r,s}_{i,j}\bigr\}$} and \smash{$\bigl\{\bar{H}^r_{i,s}\bigr\}$}. The graded algebra $\operatorname{gr}BD(n)$ is generated by \smash{$\bigl\{\bar{H}_{i,2r+1},\bar{B}_{i,r}\bigr\}$} which satisfy the following relations:
\begin{gather*}
\bigl[\bar{H}_{i,2r+1},\bar{H}_{j,2s+1}\bigr]=0,\\
\bigl[\bar{H}_{i,2r+1},\bar{B}_{j,s}\bigr]=0,\\
\bigl[\bar{B}_{i,r+1},\bar{B}_{j,s}\bigr]=\bigl[\bar{B}_{i,r},\bar{B}_{j,s+1}\bigr],\\
\bigl[\bar{B}_{i,r},\bar{B}_{j,s}\bigr]=0\qquad \text{if }a^A_{i,j}=0,\\
\bigl[\bar{B}_{i,r},\bigl[\bar{B}_{i,s},\bar{B}_{j,u}\bigr]\bigr]=0\qquad \text{if }a^A_{i,j}=-1.
\end{gather*}
Let $L_+$ be the subalgebra generated by $\bigl\{\bar{B}_{i,r}\bigr\}$.
We can decompose $\operatorname{gr}D(n)$ into the tensor product of $\mathbb{C}[H_{i,2r+1}]$ and $L_+$. By Theorem~\ref{plus}, we have a surjective homomorphism from $A(n)_+$ to $L_+$.
By definition, $\bar{E}^{r,s}_{i,j}$ corresponds to $E_{i,j}u^rv^s$, $H^r_{i,s}$ corresponds to $(E_{i,i}-E_{i+1,i+1})u^rv^s+\allowbreak {(r-1)u^rv^sdu}$ and $H^r_{0,s}$ corresponds to $(E_{n,n}-E_{1,1})u^rv^s+u^{r-1}v^sdu$. Since the set of these elements is a basis of $\bigoplus_{s\geq0}\bigoplus_{\alpha\in\Delta^+}\sum_{k=1}^{\dim \mathfrak{g}_\alpha}x_\alpha^k v^s\oplus\Omega_+$, we find that the set of ordered polynomials derived from \smash{$\bigl\{H_{k,2r+1},E_{i,j}^{r,s}\bigr\}$} spans $BD(n)$. Thus, $\phi$ becomes injective.
\end{proof}


\begin{thebibliography}{99}
\footnotesize\itemsep=0pt

\bibitem{BR1}
Belliard S., Regelskis V., Drinfeld~{J} presentation of twisted {Y}angians,
 \href{https://doi.org/10.3842/SIGMA.2017.011}{\textit{SIGMA}} \textbf{13}
 (2017), 011, 35~pages,
 \href{https://arxiv.org/abs/1401.2143}{arXiv:1401.2143}.

\bibitem{Bro}
Brown J., Twisted {Y}angians and finite {$W$}-algebras,
 \href{https://doi.org/10.1007/s00031-008-9041-x}{\textit{Transform. Groups}}
 \textbf{14} (2009), 87--114,
 \href{https://arxiv.org/abs/0710.2918}{arXiv:0710.2918}.

\bibitem{BK}
Brundan J., Kleshchev A., Shifted {Y}angians and finite {$W$}-algebras,
 \href{https://doi.org/10.1016/j.aim.2004.11.004}{\textit{Adv. Math.}}
 \textbf{200} (2006), 136--195,
 \href{https://arxiv.org/abs/math/0407012}{arXiv:math/0407012}.

\bibitem{DKV}
De~Sole A., Kac V.G., Valeri D., A~{L}ax type operator for quantum finite
 {$W$}-algebras,
 \href{https://doi.org/10.1007/s00029-018-0439-6}{\textit{Selecta
 Math.~(N.S.)}} \textbf{24} (2018), 4617--4657,
 \href{https://arxiv.org/abs/1707.03669}{arXiv:1707.03669}.

\bibitem{D1}
Drinfeld V.G., Hopf algebras and the quantum {Y}ang--{B}axter equation,
 \textit{Soviet Math. Dokl.} \textbf{32} (1985), 256--258.

\bibitem{D2}
Drinfeld V.G., A~new realization of {Y}angians and of quantum affine algebras,
 \textit{Soviet Math. Dokl.} \textbf{36} (1988), 212--216.

\bibitem{Gu1}
Guay N., Affine {Y}angians and deformed double current algebras in type~{$A$},
 \href{https://doi.org/10.1016/j.aim.2006.08.007}{\textit{Adv. Math.}}
 \textbf{211} (2007), 436--484.

\bibitem{GNW}
Guay N., Nakajima H., Wendlandt C., Coproduct for {Y}angians of affine
 {K}ac--{M}oody algebras,
 \href{https://doi.org/10.1016/j.aim.2018.09.013}{\textit{Adv. Math.}}
 \textbf{338} (2018), 865--911,
 \href{https://arxiv.org/abs/1701.05288}{arXiv:1701.05288}.

\bibitem{HU}
Harako S., Ueda M., A~new presentation of the twisted {Y}angian of type~{$D$},
 \href{https://arxiv.org/abs/2507.00350}{arXiv:2507.00350}.

\bibitem{K}
Kac V.G., Infinite-dimensional {L}ie algebras, 3rd ed.,
 \href{https://doi.org/10.1017/CBO9780511626234}{Cambridge University Press},
 Cambridge, 1990.

\bibitem{Kassel}
Kassel C., K\"ahler differentials and coverings of complex simple {L}ie
 algebras extended over a commutative algebra,
 \href{https://doi.org/10.1016/0022-4049(84)90040-9}{\textit{J.~Pure Appl.
 Algebra}} \textbf{34} (1984), 265--275.

\bibitem{L}
Lu K., Minimalistic presentation and coideal structure of twisted {Y}angians,
 \href{https://doi.org/10.1007/s00220-026-05615-3}{\textit{Comm. Math. Phys.}}
 \textbf{407} (2026), 99, 38~pages,
 \href{https://arxiv.org/abs/2511.07136}{arXiv:2511.07136}.

\bibitem{LPTTW}
Lu K., Peng Y.N., Tappeiner L., Topley L., Wang W., Shifted twisted {Y}angians
 and finite {$W$}-algebras of classical type,
 \href{https://arxiv.org/abs/2505.03316}{arXiv:2505.03316}.

\bibitem{LWZ0}
Lu K., Wang W., Zhang W., A~{D}rinfeld type presentation of twisted {Y}angians,
 \href{https://doi.org/10.1090/ert/702}{\textit{Represent. Theory}}
 \textbf{29} (2025), 838--870,
 \href{https://arxiv.org/abs/2308.12254}{arXiv:2308.12254}.

\bibitem{LWZ}
Lu K., Wang W., Zhang W., Affine {$\imath$}quantum groups and twisted
 {Y}angians in {D}rinfeld presentations,
 \href{https://doi.org/10.1007/s00220-025-05263-z}{\textit{Comm. Math. Phys.}}
 \textbf{406} (2025), 98, 36~pages,
 \href{https://arxiv.org/abs/2406.05067}{arXiv:2406.05067}.

\bibitem{LWZ1}
Lu K., Zhang W., A~{D}rinfeld type presentation of twisted {Y}angians of
 quasi-split type,
 \href{https://doi.org/10.1142/S021919972650001X}{\textit{Commun. Contemp.
 Math.}}, {t}o appear,
 \href{https://arxiv.org/abs/2408.06981}{arXiv:2408.06981}.

\bibitem{MNT}
Matsuo A., Nagatomo K., Tsuchiya A., Quasi-finite algebras graded by
 {H}amiltonian and vertex operator algebras, in Moonshine: the First Quarter
 Century and Beyond, \textit{London Math. Soc. Lecture Note Ser.}, Vol.~372,
 \href{https://doi.org/10.1017/CBO9780511730054.015}{Cambridge University
 Press}, Cambridge, 2010, 282--329,
 \href{https://arxiv.org/abs/math/0505071}{arXiv:math/0505071}.

\bibitem{MRY}
Moody R.V., Rao S.E., Yokonuma T., Toroidal {L}ie algebras and vertex
 representations, \href{https://doi.org/10.1007/BF00147350}{\textit{Geom.
 Dedicata}} \textbf{35} (1990), 283--307.

\bibitem{O}
Olshanskii G.I., Twisted {Y}angians and infinite-dimensional classical {L}ie
 algebras, in Quantum Groups ({L}eningrad, 1990), \textit{Lecture Notes in
 Math.}, Vol.~1510, \href{https://doi.org/10.1007/BFb0101183}{Springer},
 Berlin, 1992, 104--119.

\bibitem{U77}
Ueda M., Twisted affine {Y}angian and rectangular {$W$}-algebra of type~{$D$},
 \href{https://doi.org/10.1007/s10468-025-10348-3}{\textit{Algebr. Represent.
 Theory}} \textbf{28} (2025), 1195--1224,
 \href{https://arxiv.org/abs/2107.09999}{arXiv:2107.09999}.

\bibitem{U78}
Ueda M., Suggestion of a~definition of the twisted affine Yangian of
 type~{$D$}, \href{https://arxiv.org/abs/2512.06340}{arXiv:2512.06340}.

\end{thebibliography} 
\end{document}
